% !TeX spellcheck = ru_RU
% !TEX root = vkr.tex

\section{Обзор существующих решений}
\label{sec:relatedworks}

\subsection{Инструменты для морфометрии и сегментации изображений}

Задача получения морфометрических данных с изображений биологических объектов решается несколькими классами инструментов.

\textbf{ImageJ / FIJI}~\cite{Schindelin2012} является стандартом де-факто в биологических лабораториях.
Пакет предоставляет широкий набор инструментов для ручного обведения контуров и измерений,
а также плагины для полуавтоматической пороговой сегментации (Threshold, Wand tool, Trainable Weka).
Применительно к рассматриваемой задаче у FIJI выявлен ряд ограничений.
Во-первых, пороговая сегментация требует ручной настройки порогов под каждую серию снимков
и не справляется с прозрачными объектами и неоднородным фоном~--- типичными для препаратов листьев Sphagnum.
Плагин Trainable Weka Segmentation позволяет использовать машинное обучение,
однако требует ручной разметки обучающих областей для каждого нового типа объектов.
Во-вторых, рабочий процесс в FIJI фрагментирован: сегментация, коррекция, измерения
и экспорт выполняются разными инструментами, что требует экспертных знаний
и увеличивает число ручных операций.
В-третьих, FIJI не предоставляет встроенного инструмента для профильного анализа вдоль
главной оси объекта с автоматическим определением направления оси~---
поперечные сечения исследователь расставляет вручную, что вносит субъективность в измерения.

\textbf{tpsDIG2}~\cite{Rohlf2021}~--- специализированный инструмент для геометрической морфометрии,
позволяющий расставлять landmarks и строить контуры (outlines).
Требует хорошо контрастированных объектов без инородных включений и неоднородных краёв.
Листья Sphagnum с рваными краями и прозрачной текстурой плохо поддаются автоматической «обрисовке».

\textbf{PlantCV}~\cite{Fahlgren2015}~--- открытая Python-библиотека для анализа изображений растений.
Поддерживает пакетную обработку и разнообразные морфометрические функции (скелетизацию, измерение площади,
периметра, длины). Не имеет графического интерфейса, требует написания скриптов,
опирается на пороговые методы и морфологические операции без нейросетевой сегментации.

\textbf{Labelme / CVAT / Label Studio}~--- инструменты ручной разметки изображений,
позволяющие строить полигональные аннотации. Предназначены для создания обучающих наборов данных,
а не для морфометрии; не выполняют измерений.

\subsection{Сравнительный анализ}

Для сравнения выбраны критерии, непосредственно влияющие на трудоёмкость
морфометрической обработки серий гербарных снимков.

\begin{table}[H]
\centering
\caption{Сравнение существующих инструментов по ключевым критериям}
\label{tab:related}
\resizebox{\textwidth}{!}{%
\begin{tabular}{lcccccccc}
\hline
\textbf{Инструмент}
  & \textbf{\shortstack{ИИ-сегм.\\без обуч.\\данных}}
  & \textbf{\shortstack{Работа на\\сложных\\объектах}}
  & \textbf{\shortstack{Единый\\pipeline}}
  & \textbf{\shortstack{Ручная\\коррекция}}
  & \textbf{\shortstack{Профильный\\анализ\\по оси}}
  & \textbf{\shortstack{Стандарт.\\вывод}}
  & \textbf{\shortstack{Batch}}
  & \textbf{\shortstack{Без\\GPU}} \\
\hline
ImageJ / FIJI    & --    & $\pm$ & -- & + & -- & -- & $\pm$ & + \\
tpsDIG2          & --    & --    & -- & + & -- & $\pm$ & -- & + \\
PlantCV          & --    & $\pm$ & -- & -- & -- & + & +  & + \\
Labelme / CVAT   & $\pm$ & +     & -- & + & -- & -- & $\pm$ & + \\
\textbf{Разраб. система} & \textbf{+} & \textbf{+} & \textbf{+} & \textbf{+} & \textbf{+} & \textbf{+} & \textbf{+} & \textbf{+} \\
\hline
\multicolumn{9}{l}{\small \textbf{+}~--- поддерживается; $\pm$~--- частично / требует ручной настройки; \textbf{--}~--- не поддерживается} \\
\end{tabular}}
\end{table}

Сравнение показывает, что ни одно из рассмотренных решений не покрывает
полный рабочий цикл биолога в рамках одного приложения.
Наиболее близким аналогом является FIJI, однако его возможности применительно к задаче
ограничены по нескольким ключевым направлениям.

\textit{Сегментация сложных объектов.}
Встроенные инструменты FIJI (пороговая сегментация, Wand) рассчитаны на контрастные объекты.
Листья мохообразных с прозрачной клеточной структурой и рваными краями
существенно снижают качество автоматической обрисовки, вынуждая пользователя
выполнять значительный объём ручной правки.

\textit{Единый рабочий цикл.}
В FIJI сегментация, измерения и экспорт реализованы разными плагинами
и требуют переключения между инструментами и ручного переноса данных.
В разрабатываемой системе весь цикл~--- загрузка, сегментация, коррекция,
измерения, экспорт~--- выполняется в одном окне приложения без переключения контекста.

\textit{Стандартизированный профильный анализ.}
FIJI позволяет измерять ширину листа, но направление оси и точки сечения
задаются пользователем вручную для каждого образца.
Разрабатываемая система автоматически определяет главную ось методом PCA
и вычисляет ширину в фиксированных относительных точках ($0.125, 0.25, \ldots, 0.875$),
обеспечивая воспроизводимые и сопоставимые между собой профили для любой серии изображений.

\subsection{Выбор модели сегментации}

Ключевым техническим решением при разработке системы является выбор метода сегментации.
Ниже рассмотрены три класса подходов и обоснован итоговый выбор.

\subsubsection*{Требование интерактивности и проблема отбора образцов}

Прежде чем сравнивать подходы к сегментации, необходимо учесть ключевую особенность
рабочего процесса биологов-морфометристов.
На каждом изображении присутствует несколько листьев, и не все они пригодны для измерений.
Биолог отбирает «хорошие» листья~--- целые, с неповреждёнными краями, типичной для вида формы,
расположенные изолированно и не перекрывающиеся с другими объектами.
Листья с рваными краями, артефактами препарирования или атипичной морфологией исключаются.

Этот отбор по своей природе субъективен: критерии зависят от целей конкретного исследования,
вида и индивидуального опыта исследователя.
Попытка формализовать «хорошесть» листа в виде числовых порогов (площадь, компактность контура,
степень выпуклости) даёт лишь грубое приближение и требовала бы отдельной экспериментальной
валидации на реальном материале~--- задача, выходящая за рамки настоящей работы.
Таким образом, \textit{полностью автоматический} pipeline «обнаружить все листья → сегментировать →
измерить» не отвечает требованиям заказчика: он не позволяет исключить непригодные образцы
без участия эксперта.

Отсюда следует принципиальное требование к архитектуре системы:
\textbf{эксперт должен участвовать в отборе}, а система~--- снять с него трудоёмкую часть работы.
Идеальная модель сегментации в данном контексте~--- та, в которой
\textit{сам акт выбора листа и его сегментация совмещены в одно действие}:
исследователь указывает на нужный лист, система немедленно строит маску.

\subsubsection*{Классические алгоритмы компьютерного зрения}

Классические методы~--- пороговая сегментация, морфологические операции, алгоритм GrabCut,
watershed~--- не требуют обучающих данных и работают на любом CPU.
Однако они не удовлетворяют сформулированному выше требованию по двум причинам.
Во-первых, их применимость ограничена качеством входного материала:
листья Sphagnum могут быть практически прозрачными и иметь рваные края~---
именно те свойства, при которых пороговые методы дают фрагментированные маски
с пропусками в области прозрачных клеток.
Во-вторых, они не обеспечивают интерактивного отбора: алгоритм не знает,
\textit{какой именно} объект на изображении нужен биологу.

\subsubsection*{Специализированные нейросетевые модели (U-Net и аналоги)}

Архитектуры типа U-Net~\cite{Ronneberger2015} демонстрируют высокую точность сегментации биомедицинских изображений
при наличии размеченного обучающего набора данных.
В данном случае такого набора не существует: гербарная коллекция ГБС РАН содержит
тысячи образцов нескольких видов рода \textit{Sphagnum}, ни один из которых
не был ранее сегментирован в машиночитаемом формате.
Сбор репрезентативной обучающей выборки потребовал бы значительных трудозатрат
самих биологов~--- той самой ручной работы, которую требуется автоматизировать.
Помимо этого, обученная модель сегментирует все объекты подходящего класса на изображении,
не предоставляя механизма для интерактивного указания конкретного листа.

\subsubsection*{Foundation-модели сегментации: SAM и MobileSAM}

Segment Anything Model (SAM)~\cite{Kirillov2023}~--- foundation-модель сегментации от Meta AI,
обученная на датасете из более чем 1~млрд масок.
Ключевое свойство SAM~--- способность сегментировать произвольный объект
по интерактивным подсказкам (точки, ограничивающий прямоугольник) \textit{без дообучения на конкретном домене}.

SAM органично решает проблему отбора образцов: один клик по телу листа одновременно
является и выбором нужного объекта экспертом, и достаточной подсказкой для построения маски.
Биолог просто не кликает на «плохие» листья~--- и они автоматически исключаются из анализа.
Такой интерфейс не требует отдельного шага отбора и не накладывает на исследователя
никакой дополнительной когнитивной нагрузки по сравнению с тем, что он и так делает мысленно.

Базовая версия SAM имеет существенное ограничение для настольного приложения:
image encoder на основе ViT-H содержит 632~М параметров и требует GPU для приемлемой скорости работы.

\textbf{MobileSAM}~\cite{Zhang2023} решает эту проблему: тяжёлый encoder заменён
на компактный ViT-T ($\approx$5.7~М параметров), обученный методом knowledge distillation
на выходах оригинального SAM encoder.
Размер checkpoint~--- 39~МБ против 2.6~ГБ у ViT-H SAM.
Авторы показывают, что на задачах prompt-based сегментации качество MobileSAM
сопоставимо с оригиналом~\cite{Zhang2023}, при этом на CPU среднего рабочего компьютера
модель генерирует маску за 1--3~с~--- приемлемо для интерактивной работы.

\subsubsection*{Итог}

MobileSAM выбран в качестве модели сегментации по совокупности факторов:
\begin{itemize}
    \item не требует обучающих данных и обобщается на новые классы объектов;
    \item интерактивный prompt-based интерфейс естественно совмещает отбор образцов и сегментацию в одно действие;
    \item справляется с морфологически сложными объектами (прозрачные клетки, рваные края);
    \item работает на CPU, не требует GPU;
    \item умещается в дистрибутив настольного приложения (39~МБ).
\end{itemize}
SAM и MobileSAM не имеют собственного GUI и средств морфометрии~---
именно их интеграция в единый pipeline с интерактивной коррекцией и параметрическими измерениями
и составляет основной вклад разрабатываемой системы.
